\documentclass[a4paper,12pt]{article}

\usepackage[T2A]{fontenc}
\usepackage[utf8]{inputenc}
\usepackage[russian]{babel}
\usepackage{amsmath,amssymb}
\usepackage{helvet}
\renewcommand{\familydefault}{\sfdefault}
\usepackage{icomma}
\usepackage[a4paper,margin=22mm]{geometry}
\usepackage{microtype}
\usepackage{tikz}
\usetikzlibrary{calc,arrows.meta,positioning,shapes.geometric}
\usepackage{tabularx,booktabs}
\usepackage{enumitem}
\usepackage{hyperref}
\usepackage{parskip}
\usepackage{fancyhdr}
\usepackage{setspace}
\usepackage{xcolor}

\definecolor{accent}{HTML}{008080}
\definecolor{darkaccent}{HTML}{004D4D}
\definecolor{textgray}{HTML}{333333}
\definecolor{softgray}{HTML}{6E7474}
\definecolor{lightaccent}{HTML}{EAF4F4}

\color{textgray}
\onehalfspacing
\setlength{\parindent}{0pt}
\setlength{\headheight}{25pt}

\hypersetup{
    colorlinks=true,
    linkcolor=darkaccent,
    urlcolor=darkaccent,
    pdfauthor={Лёвин Артём Александрович},
    pdftitle={Чек-лист: уравнения, корни и равносильные уравнения}
}

\pagestyle{fancy}
\fancyhf{}
\fancyhead[L]{\small\color{softgray}Кирилл Зиновьев}
\fancyhead[R]{\small\color{softgray}7 класс}
\fancyfoot[C]{\small\color{softgray}\thepage}
\renewcommand{\headrulewidth}{0.3pt}
\renewcommand{\headrule}{%
  \hbox to\headwidth{\color{accent}\leaders\hrule height \headrulewidth\hfill}%
}

\setlist[itemize]{leftmargin=1.6em,itemsep=0.35em,topsep=0.4em}
\setlist[enumerate]{leftmargin=1.8em,itemsep=0.65em,topsep=0.5em}

\newcommand{\sectiontitle}[1]{%
    \vspace{0.35em}%
    {\Large\bfseries\color{darkaccent}#1\par}%
    \vspace{0.2em}%
    {\color{accent}\rule{\linewidth}{0.7pt}}%
    \vspace{0.6em}%
}

\newcommand{\subsectiontitle}[1]{%
    \vspace{0.55em}%
    {\large\bfseries\color{darkaccent}#1\par}%
    \vspace{0.25em}%
}

\newcommand{\important}[1]{%
    \par\vspace{0.3em}%
    \noindent
    \begin{minipage}{\linewidth}
        {\color{accent}\rule{2.2pt}{\baselineskip}}\hspace{0.6em}%
        \begin{minipage}[t]{0.94\linewidth}
            #1
        \end{minipage}
    \end{minipage}%
    \par\vspace{0.3em}%
}

\tikzset{
    flowstart/.style={
        ellipse,
        draw=darkaccent,
        line width=0.9pt,
        align=center,
        minimum width=4.6cm,
        minimum height=1cm,
        fill=lightaccent,
        font=\small
    },
    flowaction/.style={
        rectangle,
        rounded corners=3pt,
        draw=accent,
        line width=0.9pt,
        align=center,
        text width=10.2cm,
        minimum height=1.1cm,
        inner sep=7pt,
        font=\small
    },
    flowdecision/.style={
        diamond,
        aspect=2.45,
        draw=darkaccent,
        line width=0.9pt,
        align=center,
        text width=3.7cm,
        inner sep=2pt,
        font=\small
    },
    flowend/.style={
        ellipse,
        draw=darkaccent,
        line width=0.9pt,
        align=center,
        text width=4.2cm,
        minimum height=1cm,
        fill=lightaccent,
        font=\small
    },
    flowarrow/.style={
        -{Stealth[length=3mm,width=2mm]},
        line width=0.9pt,
        draw=darkaccent
    }
}

\begin{document}

\begin{titlepage}
\thispagestyle{empty}

\begin{center}
\vspace*{2.0cm}

{\color{accent}\rule{0.68\linewidth}{1.1pt}}

\vspace{1.0cm}

{\Huge\bfseries Чек-лист\par}

\vspace{0.7cm}

{\LARGE\bfseries\color{darkaccent}
Уравнения, корни\\[0.15cm]
и равносильные уравнения
\par}

\vspace{0.6cm}

{\large
Раскрытие скобок и приведение подобных слагаемых
\par}

\vspace{1.6cm}

{\large Кирилл Зиновьев\par}
\vspace{0.25cm}
{\large 7 класс\par}

\vfill

{\large
Лёвин Артём Александрович\par
эксклюзивно для Кирилла Зиновьева
}

\vspace{0.7cm}

{\large \today}

\vspace{2.1cm}
\end{center}

\begin{tikzpicture}[remember picture,overlay]
    \draw[
        accent,
        line width=1.2pt,
        opacity=0.55
    ]
    ($(current page.south west)+(1.7cm,1.6cm)$)
    .. controls
    ($(current page.south west)+(6.5cm,2.7cm)$)
    and
    ($(current page.south east)+(-6.8cm,0.8cm)$)
    ..
    ($(current page.south east)+(-1.7cm,1.7cm)$);

    \draw[
        darkaccent,
        line width=0.5pt,
        opacity=0.35
    ]
    ($(current page.south west)+(2.3cm,1.15cm)$)
    .. controls
    ($(current page.south west)+(7.7cm,2.0cm)$)
    and
    ($(current page.south east)+(-6.0cm,1.05cm)$)
    ..
    ($(current page.south east)+(-2.5cm,1.25cm)$);
\end{tikzpicture}

\end{titlepage}

\sectiontitle{1. Что нужно уметь}

После повторения материала Вы должны уметь:

\begin{itemize}
    \item раскрывать скобки при умножении выражения на число;
    \item правильно определять знаки после раскрытия скобок;
    \item находить и приводить подобные слагаемые;
    \item понимать, что называется корнем уравнения;
    \item проверять подстановкой, является ли число корнем;
    \item находить корни простейших уравнений с квадратом и модулем;
    \item сравнивать множества корней и распознавать равносильные уравнения.
\end{itemize}

\sectiontitle{2. Раскрытие скобок}

Если перед скобками стоит множитель, его нужно умножить на каждое
слагаемое внутри скобок:

\[
a(b+c)=ab+ac,
\qquad
a(b-c)=ab-ac.
\]

Например,

\[
2(x+5)=2x+10.
\]

Если множитель отрицательный, правило остаётся тем же:

\[
-2(x+5)=-2x-10,
\]

\[
-2(-x+5)=2x-10.
\]

\important{
Знак каждого нового слагаемого определяется обычным правилом знаков
при умножении. Особенно внимательно следите за отрицательным множителем.
}

Полезно сравнить четыре случая:

\[
-10(2-3x)=-20+30x,
\]
\[
-10(-2-3x)=20+30x,
\]
\[
-10(2+3x)=-20-30x,
\]
\[
-10(-2+3x)=20-30x.
\]

\subsectiontitle{Минус перед скобками}

Минус перед скобками означает умножение всей скобки на \(-1\):

\[
-(a-b)=-a+b.
\]

Поэтому

\[
-(a+b)=-a-b,
\qquad
-(a-b)=-a+b.
\]

\sectiontitle{3. Подобные слагаемые}

\textbf{Подобные слагаемые} имеют одинаковую буквенную часть.

Например, в выражении

\[
3x+5x-4y+2y
\]

подобны пары

\[
3x \text{ и } 5x,
\qquad
-4y \text{ и } 2y.
\]

При приведении подобных складываются их числовые коэффициенты,
а буквенная часть сохраняется:

\[
3x+5x=8x,
\qquad
-4y+2y=-2y.
\]

Следовательно,

\[
3x+5x-4y+2y=8x-2y.
\]

Ещё один пример:

\[
-2k+4t-7k-t=-9k+3t.
\]

\important{
Знак относится к слагаемому целиком. Например, коэффициент слагаемого
\(-7k\) равен \(-7\), а коэффициент слагаемого \(+4t\) равен \(4\).
}

\sectiontitle{4. Сначала скобки, затем подобные}

Типовой порядок действий:

\begin{enumerate}
    \item раскрыть скобки;
    \item проверить получившиеся знаки;
    \item найти подобные слагаемые;
    \item сложить их коэффициенты.
\end{enumerate}

Пример:

\[
2(x-1)-3(5-2x).
\]

Раскрываем скобки:

\[
2x-2-15+6x.
\]

Приводим подобные:

\[
2x+6x-2-15=8x-17.
\]

Итог:

\[
\boxed{8x-17}.
\]

\sectiontitle{5. Корень уравнения}

\textbf{Корень уравнения} --- это значение переменной, при котором
левая и правая части уравнения имеют одинаковые значения.

Например,

\[
x+3=5.
\]

При \(x=2\):

\[
2+3=5.
\]

Получилось верное равенство, поэтому

\[
\boxed{x=2}
\]

является корнем уравнения.

\important{
Чтобы проверить предполагаемый корень, нужно подставить его вместо
\textbf{каждого} вхождения переменной и отдельно вычислить левую и правую части.
}

\subsectiontitle{Проверка подстановкой}

Рассмотрим уравнение

\[
3x^2=5x+12.
\]

Проверим \(x=1\):

\[
3\cdot1^2=3,
\qquad
5\cdot1+12=17.
\]

Так как

\[
3\ne17,
\]

число \(1\) корнем не является.

Проверим \(x=3\):

\[
3\cdot3^2=27,
\qquad
5\cdot3+12=27.
\]

Получили верное равенство:

\[
27=27.
\]

Следовательно,

\[
\boxed{x=3}
\]

является корнем уравнения.

\important{
Проверка одного числа отвечает только на вопрос, является ли корнем
именно это число. Такая проверка сама по себе не находит все остальные корни.
}

\clearpage
\thispagestyle{plain}

\begin{center}
{\Large\bfseries\color{darkaccent}
Алгоритм: как проверить, является ли число корнем уравнения}
\end{center}

\vspace{0.8cm}

\begin{center}
\begin{tikzpicture}[node distance=0.9cm]

\node[flowstart] (start)
{Дано уравнение и число \(a\)};

\node[flowaction, below=of start] (sub)
{Подставьте \(x=a\) вместо каждого \(x\) в исходном уравнении};

\node[flowaction, below=of sub] (left)
{Вычислите значение левой части};

\node[flowaction, below=of left] (right)
{Вычислите значение правой части};

\node[flowdecision, below=1.15cm of right] (compare)
{Полученные значения равны?};

\node[flowend, below=1.55cm of compare, xshift=-3.5cm] (yes)
{\(a\) --- корень уравнения};

\node[flowend, below=1.55cm of compare, xshift=3.5cm] (no)
{\(a\) --- не корень уравнения};

\draw[flowarrow] (start) -- (sub);
\draw[flowarrow] (sub) -- (left);
\draw[flowarrow] (left) -- (right);
\draw[flowarrow] (right) -- (compare);

\draw[flowarrow]
(compare.south west)
-- node[pos=0.35,left,font=\small\bfseries]{да}
(yes.north);

\draw[flowarrow]
(compare.south east)
-- node[pos=0.35,right,font=\small\bfseries]{нет}
(no.north);

\end{tikzpicture}
\end{center}

\clearpage

\sectiontitle{6. Уравнения с модулем}

Модуль числа показывает его расстояние от нуля. Поэтому

\[
|-4|=4,
\qquad
|4|=4.
\]

Уравнение

\[
|x|=4
\]

имеет два корня:

\[
\boxed{x=-4},
\qquad
\boxed{x=4}.
\]

В общем случае для уравнения

\[
|x|=a
\]

получаем:

\begin{itemize}
    \item если \(a>0\), то \(x=a\) или \(x=-a\);
    \item если \(a=0\), то \(x=0\);
    \item если \(a<0\), корней нет, потому что \(|x|\ge0\).
\end{itemize}

\sectiontitle{7. Произведение, равное нулю}

Если произведение равно нулю, то хотя бы один из множителей равен нулю:

\[
AB=0
\quad\Longrightarrow\quad
A=0 \text{ или } B=0.
\]

Например,

\[
(x-5)(x+5)=0.
\]

Отсюда

\[
x-5=0
\qquad\text{или}\qquad
x+5=0,
\]

поэтому

\[
\boxed{x=5}
\qquad\text{или}\qquad
\boxed{x=-5}.
\]

\sectiontitle{8. Равносильные уравнения}

Два уравнения называются \textbf{равносильными}, если множества их
корней полностью совпадают.

Рассмотрим:

\[
\text{A)}\quad x^2=25,
\]
\[
\text{B)}\quad x-5=0,
\]
\[
\text{C)}\quad (x-5)(x+5)=0.
\]

Для уравнения A:

\[
x=5 \quad\text{или}\quad x=-5,
\]

поэтому множество корней равно

\[
\{-5,5\}.
\]

Для уравнения B:

\[
x=5,
\]

поэтому множество корней равно

\[
\{5\}.
\]

Для уравнения C:

\[
x=5 \quad\text{или}\quad x=-5,
\]

поэтому множество корней равно

\[
\{-5,5\}.
\]

Следовательно,

\[
\boxed{\text{A и C равносильны}.}
\]

\important{
Для равносильности недостаточно одного общего корня. Должны полностью
совпадать все корни рассматриваемых уравнений.
}

Ещё один пример:

\[
\text{A)}\quad x=6,
\qquad
\text{B)}\quad x^2=36,
\qquad
\text{C)}\quad (x-6)(x+6)=0.
\]

Множества корней:

\[
A:\ \{6\},
\qquad
B:\ \{-6,6\},
\qquad
C:\ \{-6,6\}.
\]

Значит, уравнения B и C равносильны.

\sectiontitle{9. Типичные ошибки}

\begin{itemize}
    \item Внешний множитель умножают только на первое слагаемое в скобках.

    \item При отрицательном множителе ошибаются со знаком одного из слагаемых.

    \item Складывают слагаемые с разной буквенной частью. Например,
    выражение \(3x+2y\) привести к одному слагаемому нельзя.

    \item При проверке корня подставляют число только вместо одного из нескольких
    вхождений переменной.

    \item Делают вывод о корне до того, как вычислены и сравнены обе части уравнения.

    \item В уравнении \(x^2=a\), где \(a>0\), забывают отрицательный корень.

    \item В уравнении \(|x|=a\), где \(a>0\), записывают только один корень.

    \item Считают уравнения равносильными, если у них есть хотя бы один общий корень.
\end{itemize}

\sectiontitle{10. Краткая памятка}

\begin{enumerate}
    \item
    \[
    a(b+c)=ab+ac.
    \]

    \item
    \[
    -(a-b)=-a+b.
    \]

    \item Подобные слагаемые имеют одинаковую буквенную часть.

    \item Корень уравнения --- значение переменной, при котором левая и правая
    части равны.

    \item Проверка корня выполняется подстановкой.

    \item
    \[
    |x|=a,\quad a>0
    \quad\Longrightarrow\quad
    x=\pm a.
    \]

    \item
    \[
    AB=0
    \quad\Longrightarrow\quad
    A=0 \text{ или } B=0.
    \]

    \item Равносильные уравнения имеют одинаковые множества корней.
\end{enumerate}

\clearpage
\sectiontitle{Тренировочный блок. Путь Б}

\textbf{Период: 24.09.2026--25.09.2026.}

Выполняйте задания письменно. В заданиях на проверку корня показывайте
подстановку и сравнение левой и правой частей.

\subsectiontitle{24.09.2026}

\begin{enumerate}[label=\textbf{\arabic*.}]
    \item Раскройте скобки и приведите подобные:
    \[
    -3(2x-5)+4x-7.
    \]

    \item Определите, является ли число \(x=2\) корнем уравнения
    \[
    2x^2+3=5x+1.
    \]

    \item Найдите все корни уравнения:
    \[
    |x|=7.
    \]
\end{enumerate}

\vspace{0.8cm}

\subsectiontitle{25.09.2026}

\begin{enumerate}[label=\textbf{\arabic*.}]
    \item Раскройте скобки и приведите подобные:
    \[
    -4(3-2x)-5(x+1).
    \]

    \item Среди чисел
    \[
    -2,\quad -1,\quad 1,\quad 2
    \]
    укажите все корни уравнения
    \[
    3x^2=5x+2.
    \]

    \item Определите, какие из уравнений равносильны:
    \[
    \text{A)}\quad x^2=49,
    \]
    \[
    \text{B)}\quad x-7=0,
    \]
    \[
    \text{C)}\quad (x-7)(x+7)=0.
    \]
\end{enumerate}

\clearpage
\sectiontitle{Ответы}

\renewcommand{\arraystretch}{1.4}

\begin{table}[ht]
\centering
\begin{tabularx}{\linewidth}{
    @{}
    >{\bfseries}p{2.8cm}
    >{\centering\arraybackslash}p{1.2cm}
    X
    @{}
}
\toprule
Дата & № & Ответ \\
\midrule
24.09.2026 & 1 & \(\displaystyle -2x+8\) \\
24.09.2026 & 2 & Да, \(x=2\) является корнем. \\
24.09.2026 & 3 & \(\displaystyle x=-7,\;x=7\) \\
\midrule
25.09.2026 & 1 & \(\displaystyle 3x-17\) \\
25.09.2026 & 2 & Из перечисленных чисел корнем является только \(x=2\). \\
25.09.2026 & 3 & Уравнения A и C; их множество корней \(\{-7,7\}\). \\
\bottomrule
\end{tabularx}
\end{table}

\vfill

\begin{center}
{\color{accent}\rule{0.55\linewidth}{0.7pt}}

\vspace{0.5em}

{\small\color{softgray}Лёвин Артём Александрович}
\end{center}

\end{document}